This chapter defines the \emph{technology frame}: a news frame attached to a named energy technology and measured through concrete frame elements. Section~\ref{subsec:framing} contrasts Entman's definition with the codebook's more permissive elements and traces the nine frames to their sources. Section~\ref{subsec:frame_elements} explains why seventeen presence questions replace holistic frame labels. Section~\ref{subsec:relation} defines the record tying a technology to its frames and their evaluation, shown and scored in Figure~\ref{fig:anatomy}.

\subsection{Framing}\label{subsec:framing}

Following \textcite[p.~52]{entman1993framing}, to frame is \enquote{to select some aspects of a perceived reality and make them more salient in a communicating text}. Such a selection can define a problem, interpret causes, evaluate morally or recommend a treatment \parencite[p.~52]{entman1993framing}. The codebook requires neither such a function nor a complete argument. Most elements are marked when a paragraph refers to their subject for a technology, without a stated mechanism. Only bare naming and purely procedural or factual reporting stay unframed, a boundary that is not sharp (Section~\ref{subsec:limitations_future}).

The evaluation records the direction in which a frame presents a technology. This operational codebook dimension is close to but not identical with Entman's moral evaluation. \textcite[p.~57]{entman1993framing} criticises coders who \enquote{simply tote up} positive and negative messages without measuring salience. Tying each evaluation to a frame answers this only in part: presence is counted, not salience, and only in the text.

The frame dimension operationalises issue framing in Entman's lineage: which aspects of an issue a naturally occurring text makes salient \parencite{card2015mediaframes}. Valence framing, a second tradition, presents the same information positively or negatively, and its attribute variant frames one characteristic of an object, such as beef described by its lean rather than its fat share \parencite[pp.~158--159]{levin1998allframes}. The evaluation dimension does not measure this. Like tone or aspect polarity, it records the direction of what a paragraph asserts within a codebook-defined frame.

Entman lists no issue-specific frames. The topic description supplies five: costs, jobs and the economy, climate and emissions, energy security, and reliability and system integration. Before calibration, these were related to the Policy Frames Codebook's issue-general dimensions \parencite{boydstun2020policyframes,card2015mediaframes} and the energy context labels of \textcite{toetzke2022energydiscourse} (Table~\ref{tab:frame_crosswalk}). Technical feasibility, siting and regulation were added, without a recorded mapping, after the first calibration session, in which 9 of 20 paragraphs fitted no element. A ninth, residual frame follows the Policy Frames Codebook and makes inventory misfit measurable.

\subsection{Frame Elements Instead of Frame Labels}\label{subsec:frame_elements}

Element-based coding is described as less subjective \parencite[p.~3]{ali2025hydrogenframing}, and \textcite[p.~264]{matthes2008contentanalysis} argue that single frame elements are more reliable than abstract, holistic frames. Yet their elements are variables of Entman's four functions, clustered into frames across texts \parencite[pp.~263--264]{matthes2008contentanalysis}. The argument therefore motivates decomposition but establishes no advantage for this study's mostly topical indicators. Two structurally close studies derive frames from binary indicators. \textcite[p.~7]{ali2025hydrogenframing} code sub-variables of Entman's four elements, and \textcite[p.~987]{liu2023strategyframing} mark a frame present when any of its items occurs.

As presence questions, elements let humans and models answer the same questions, on a coding sheet or in a prompt. Per-element agreement also keeps pooled figures from hiding weak elements. \textcite{matthes2008contentanalysis} and \textcite{ali2025hydrogenframing} derive frames from data. This thesis instead fixes the mapping of seventeen elements to nine frames beforehand as its own operationalisation: a frame is present if any of its elements is (Table~\ref{tab:codebook_elements}). The coder's two codings measure temporal stability, but whether derived frames agree with holistic frame judgements remains untested.

\subsection{The Technology--Frame--Evaluation Relation}\label{subsec:relation}\label{subsec:domain_background}

The pipeline and coder produce one \emph{technology record}, or row, per technology a paragraph makes a statement about. The closed inventory comprises renewable energy in general, solar, wind, hydrogen, other renewables, energy storage, and nuclear energy as a low-carbon comparator. Wind siting and, following the topic description, the hydrogen production pathway \parencite{sadatrazavi2024hydrogen} are exported as subtypes but validated only for execution stability, not against the human reference. A comparison target records what the technology is compared with. Elements that \emph{fire}, or are marked present, determine the record's frames by the fixed disjunction. Each frame carries one evaluation (Figure~\ref{fig:anatomy}).

% 2026-11-09, second slimming pass (G-54): caption syntax tightened; no fact removed.
% 2026-11-09: caption shortened (consensus F-08; the author's note: even shorter)
% Figure 2 (Chapter 2, reused for scoring in Chapter 5) -- anatomy of a technology frame and the scoring of a model
% output against it. Drawn by hand in TikZ in the thesis figure style (preamble); restructuring plan v3, Fig 2.
% THE PARAGRAPH IS INVENTED. It is not corpus text and paraphrases none: checked on 2026-11-05 against all 1,047
% paragraphs and headlines of tools/annotation/sheets/ALL_paragraphs.csv (lower-cased word 5-grams): no shared
% 5-word sequence (script and counts only, no corpus text printed). Do not replace it with a corpus paragraph.
% Element names, frames, the disjunctive derivation, the evaluation values and the comparison-target value
% "natural gas" (natural_gas) are those of tools/annotation/codebook/codebook-v1.0.yaml. The coding in panel a is the
% author of the figure's reading of the invented text under the codebook, for illustration only; the model output in
% panel b is hypothetical. The scoring follows tools/pipeline/evaluate.py: element tuples (paragraph x technology x
% element) scored as sets; evaluation score A on the reference (row, frame) pairs the model also predicted, score B on
% (paragraph, technology, frame, value) tuples; comparison target per row, score A on rows present in both;
% displacement on paragraph-level presence indicators (Equation eq:displacement). Arithmetic of panel b: elements TP 2,
% FP 1, FN 2 -> precision 2/3, recall 2/4, micro-F1 4/7 = 0.57; evaluation score A 2/2 = 1.00, score B TP 2, FP 1,
% FN 1 -> 4/6 = 0.67; comparison target: the model gives natural gas (wind) and none (solar), as the coder does in panel a,
% so 2 of 2 matched rows agree (shown in the footer since 2026-11-05, so that the reader can verify it).
% Colour doubles the written outcome only (missed in model blue, invented in orange, as "below" and "above" the
% reference elsewhere); every outcome is also spelled out.
% Palette 2026-11-06: outcome cells moved to the palette names (missed: figModel, invented: figAbove, matched and TP:
% figInk), the "Coder" and "Model" column heads of panel b marked figInk and figModel, the caption names the two
% colours; no value, label text or band fill changed, panel a untouched.
\begin{figure}[htbp]
\centering
\begin{tikzpicture}[
  font=\figfontsmall,
  para/.style={draw=figInk, line width=0.5pt, rectangle, rounded corners=1.5pt, fill=white, align=left,
    inner sep=4pt, text width=146mm, font=\figfont},
  band/.style={fill=figGreyC, rounded corners=2pt},
  el/.style={draw=figInk, line width=0.5pt, rectangle, rounded corners=1.5pt, fill=white, align=left,
    anchor=west, minimum height=5.6mm, text width=63mm, inner sep=2pt},
  fr/.style={draw=figInk, line width=0.8pt, rectangle, rounded corners=1.5pt, fill=white, align=center,
    minimum height=5.6mm, text width=19mm, inner sep=2pt},
  head/.style={anchor=base west, text=figGreyA, font=\figfontsmall},
  tech/.style={anchor=west, font=\figfont\bfseries},
  val/.style={anchor=west},
  arr/.style={-{Latex[length=1.6mm]}, line width=0.6pt},
  cue/.style={},
]
% ---------------------------------------------------------------- panel a: the paragraph
\node[para, anchor=north west] (p) at (0,0) {%
The new offshore wind farm will supply electricity \textbf{more cheaply than gas-fired power stations}\textsuperscript{1} and \textbf{needs no state subsidy}\textsuperscript{2}, the operator said. Residents on the coast fear that \textbf{the turbines will spoil their view of the sea}\textsuperscript{3}. Meanwhile, \textbf{planning law still blocks a solar park}\textsuperscript{4} next to the harbour.};
\node[anchor=south west, inner sep=0pt, font=\figfontsmall\itshape, text=figGreyA] at ([yshift=1.2mm]p.north west)
  {Constructed example, invented for this figure (not corpus text)};
\node[panellabel, anchor=south east] at ([xshift=-1.5mm, yshift=1.2mm]p.north west) {a};

% ---------------------------------------------------------------- panel a: the coding of that paragraph
\begin{scope}[shift={(0,-21mm)}]
  \node[head] at (0,0) {Technology};
  \node[head] at (21mm,0) {Fired elements (codebook names)};
  \node[head] at (95mm,0) {Frame};
  \node[head] at (118mm,0) {Evaluation};
  \node[head, align=left] at (135mm,0) {Comparison\\target};
  \node[head, font=\figfontsmall\itshape] at (95mm,-3.4mm) {any element fires};
  % wind row
  \fill[band] (0,-6mm) rectangle (152mm,-31mm);
  \node[tech] at (1.5mm,-18mm) {Wind};
  \node[el] (e1) at (21mm,-10mm) {\textsuperscript{1}\,Cost or price (F1a)};
  \node[el] (e2) at (21mm,-18mm) {\textsuperscript{2}\,Subsidy, levy or public funding (F1b)};
  \node[el] (e3) at (21mm,-26mm) {\textsuperscript{3}\,Landscape, visual impact and local amenity (F7b)};
  \node[fr] (f1) at (105mm,-14mm) {Costs};
  \node[fr] (f2) at (105mm,-26mm) {Siting};
  \draw[arr] (e1.east) -- ++(4mm,0) |- ([yshift=1mm]f1.west);
  \draw[arr] (e2.east) -- ++(4mm,0) |- ([yshift=-1mm]f1.west);
  \draw[arr] (e3.east) -- (f2.west);
  \node[val] at (118mm,-14mm) {positive};
  \node[val] at (118mm,-26mm) {negative};
  \node[val] at (135mm,-18mm) {natural gas};
  % solar row
  \fill[band] (0,-33mm) rectangle (152mm,-41mm);
  \node[tech] at (1.5mm,-37mm) {Solar};
  \node[el] (e4) at (21mm,-37mm) {\textsuperscript{4}\,Regulation, permitting and legal process (F8a)};
  \node[fr] (f3) at (105mm,-37mm) {Regulation};
  \draw[arr] (e4.east) -- (f3.west);
  \node[val] at (118mm,-37mm) {negative};
  \node[val] at (135mm,-37mm) {none};
\end{scope}

% ---------------------------------------------------------------- panel b: scoring a hypothetical model output
\begin{scope}[shift={(0,-71mm)}]
  \node[panellabel, anchor=south east] at (-1.5mm,1.2mm) {b};
  \node[anchor=south west, inner sep=0pt, font=\figfontsmall\itshape, text=figGreyA] at (0,1.2mm)
    {A hypothetical model output scored against the coding in panel a};
  % elements
  \node[anchor=north west, inner sep=0pt] (tb1) at (0,-1mm) {%
    \renewcommand{\arraystretch}{1.12}\setlength{\tabcolsep}{3pt}%
    \begin{tabular}{@{}l c c l@{}}
    \multicolumn{4}{@{}l}{\textbf{Frame elements}, tuples (row $\times$ element)} \\
    \toprule
    Row $\cdot$ element & \textcolor{figInk}{Coder} & \textcolor{figModel}{Model} & Scored as \\
    \midrule
    wind $\cdot$ cost or price (F1a) & 1 & 1 & \textcolor{figInk}{TP} \\
    wind $\cdot$ subsidy (F1b) & 1 & 0 & \textcolor{figModel}{FN, missed} \\
    wind $\cdot$ landscape (F7b) & 1 & 0 & \textcolor{figModel}{FN, missed} \\
    wind $\cdot$ emissions (F3a) & 0 & 1 & \textcolor{figAbove}{FP, invented} \\
    solar $\cdot$ regulation (F8a) & 1 & 1 & \textcolor{figInk}{TP} \\
    \midrule
    \multicolumn{4}{@{}l}{precision 2/3, recall 2/4, micro-F1 0.57} \\
    \bottomrule
    \end{tabular}};
  % evaluation and comparison target
  \node[anchor=north west, inner sep=0pt] (tb2) at (78mm,-1mm) {%
    \renewcommand{\arraystretch}{1.12}\setlength{\tabcolsep}{3pt}%
    \begin{tabular}{@{}l l l l@{}}
    \multicolumn{4}{@{}l}{\textbf{Evaluation}, one value per row $\times$ frame} \\
    \toprule
    Row $\cdot$ frame & \textcolor{figInk}{Coder} & \textcolor{figModel}{Model} & Scored as \\
    \midrule
    wind $\cdot$ costs & positive & positive & \textcolor{figInk}{matched, agrees} \\
    wind $\cdot$ siting & negative & -- & \textcolor{figModel}{FN, frame missed} \\
    wind $\cdot$ climate & -- & positive & \textcolor{figAbove}{FP, frame invented} \\
    solar $\cdot$ regulation & negative & negative & \textcolor{figInk}{matched, agrees} \\
    \midrule
    \multicolumn{4}{@{}l}{matched-frame accuracy 2/2 = 1.00} \\
    \multicolumn{4}{@{}l}{end-to-end F1 0.67 (TP 2, FP 1, FN 1)} \\
    \multicolumn{4}{@{}l}{comparison target, model: natural gas (wind), none (solar);} \\
    \multicolumn{4}{@{}l}{as coded, so 2 of 2 matched rows agree} \\
    \bottomrule
    \end{tabular}};
  % from tuples to a paragraph prevalence
  \node[anchor=north west, inner sep=0pt] at (0,-50mm) {\textbf{Paragraph indicator}: element present on any row};
  \node[anchor=north west, inner sep=0pt] (tb3) at (0,-54.5mm) {%
    \renewcommand{\arraystretch}{1.12}\setlength{\tabcolsep}{3pt}%
    \begin{tabular}{@{}l *{5}{>{\centering\arraybackslash}p{7mm}}@{}}
    \toprule
    & F1a & F1b & F3a & F7b & F8a \\
    \midrule
    \textcolor{figInk}{Coder}, $y_i$ & 1 & 1 & 0 & 1 & 1 \\
    \textcolor{figModel}{Model}, $\hat{y}_i$ & 1 & \textcolor{figModel}{0} & \textcolor{figAbove}{1} & \textcolor{figModel}{0} & 1 \\
    \bottomrule
    \end{tabular}};
  \node[anchor=north west, text width=72mm, align=left, inner sep=0pt] at (78mm,-54.5mm) {%
    Prevalence: the design-weighted share of sampled paragraphs with indicator 1; displacement: the model's share minus the coder's, in percentage points, as in Equation~\eqref{eq:displacement}. Here, the model's share falls for subsidy and landscape and rises for emissions.};
  \path (0,-75mm); % keeps a little space above the caption
\end{scope}
\end{tikzpicture}
\caption[Anatomy of a technology frame and scoring of a model output]{\textbf{Technology frames derive from fired elements; model output is scored at three levels.} \textbf{a},~An invented paragraph sharing no five-word sequence with the corpus. \textbf{b},~The hypothetical output misses subsidy, but cost or price keeps the costs frame present. TP, FP, FN: true positive, false positive, false negative. Blue: the model and its misses; orange: its inventions.}
\label{fig:anatomy}
\end{figure}


The evaluation is positive, negative, neutral or mixed. Its direction follows the Policy Frames Codebook's litmus test of whether a supporter or opponent would welcome the text \parencite[p.~35]{boydstun2020policyframes}, with the test moved within each frame from a policy to a technology. The result is a coding judgement, not observed audience response or outlet stance. When torn between neutral and a valenced label, the model is told to choose neutral, so unclear and neutral text can share a label. The topic description's fifth value, uncertain, becomes a separate coder-confidence rating, so uncertainty in the text is not recorded.

Each frame of a technology record is a \emph{technology frame}. Unlike the \emph{technological frames} that groups of people hold \parencite{orlikowski1994frames}, it is a property of a text, and a record can carry several or none. Technology, frame and evaluation correspond to the entity, the aspect discussed and its polarity in aspect-category sentiment analysis \parencite{pontiki2014semeval,pontiki2016semeval}. The value \emph{mixed} follows the 2014 Semantic Evaluation (SemEval) shared task's \emph{conflict} label, which forced no dominant polarity \parencite{pontiki2014semeval}.
